\preludeandfugue{Death of a Piano}
\contributor{Brynneth Nimue Brown}

% Memory, melancholy and loss

\prelude

Modern pianos have steel frames to prevent the wood from warping. My
grandmother's piano survived for about a century, in spite of only
being wood. Built for a cold home, the piano could not cope with
central heating and in its final years, became untunable. No matter
the sentimental value, a century-old piano is too large a thing to
keep when it can no longer be played. The memory of it lives on in me.

\fugue

I found the small, dead body on the stairs, pitiful and soggy from
having been chewed. The cat who had previously brought the little
rodent in had inevitably caught it after a few weeks of hide and
seek. I had been aware of the mouse and had surreptitiously been
putting out food and water. Partly to give the mouse a fighting chance
of survival. Partly in the hopes of protecting the many things a
hungry mouse might otherwise nibble.

The mouse had been living inside my grandmother's piano. I found this
out during the summer while watering the piano. We had been advised by
a tuner that a pot of water inside the body would help prevent the
wood from drying out in hotter weather. I duly took responsibility for
piano watering. When I removed the dark board from above the pedals, I
found a little stash of dried pea flakes. Food I had put out for the
mouse had been stashed here, in a furtive, hopeful way.

I wondered how often the mouse had been inside the piano while I was
playing. What sense could it possibly have made of the sound? I
worried that I might have frightened it. At the same time, I hoped
that it had somehow been able to appreciate the music as much as I
did.

The piano came into the house when my grandmother was a child. It
brought sorrow in the form of an abusive music teacher, but it must
also have brought solace. She kept it for most of her life and
persisted in playing it until her hands could no longer manage the
keys. Whatever it had cost her, she evidently considered it worth the
price. The piano outlived her marriage, and her parents, enduring to
watch over my childhood and to accept new small fingers on the
keys. It came with stories of wilder times, parties and musicians. I
grew up with the music she chose; classical pieces played with more
passion than skill, and a relentless determination to improve.

She did not really play for other people; the music held a private
place in her life. If she had ever aspired to be professional, she
never mentioned it. She was always the hobbyist where her brother was
the professional. This held true in both music and art, where he had a
career and some side hustles, and she had her quiet dedication.

I do not know what she dreamed of, or where her mind went when her
fingers were on the keys. She held no ambition on my behalf, never
spoke of success or achievement. She simply wanted me to play, and
there was a gift in that, which offered possibility, while making no
demands.

In old age, a stroke robbed her of the use of her left hand. She
fought to regain motor control. Fought relentlessly to claw back what
had been taken from her. She played, badly and persistently, ringing
defiant notes with a left hand no longer able to keep up. There were
difficult pieces she eventually abandoned but sheer bloody-mindedness
brought some material back under her fingers. She refused defeat.

When she could no longer play, she insisted that I must have the
piano. She wanted my small son to learn. For some years, the piano sat
in my front room, exiled from its home, and a little lost. There is a
stoicism in pianos. I could not bear to think about how that other
room would be without it. A room that for my entire life had existed
only for piano music. Emptiness is more than a question of unfilled
space; death can seem like an event, but it is more normally a
process.

No doubt the mouse who lived in the piano died quickly, but the
process of that death began weeks previously, in capture. A piano dies
slowly as time takes a toll on the structure of wood. Dreams are also
slow to die. Hope fades in a life that is not lived fully. The piano
stands untouched, accumulating dust. Marriages die in small doses of
everyday poison, in carelessness and casual cruelty. Homes are
abandoned, envisaged futures melt away, leaving only stains and
regrets. Sometimes it is a kindness when memory fades.

After my grandmother's death. After the death of her piano.

I remember the dark green cover of the preludes and fugues. I did not
choose them but had them thrust upon me by a teacher who had no
interest in my preferences. There are many ways to kill a dream. I
have no memory of specific notes on the stave, only the density of it
and my own annoyance. Like my grandmother before me, I played with
more passion than skill and Bach did not agree with me. I remember
frustration, the death of joy in my own inability to make this music
come alive under my hands. I can still remember the feeling of those
old keys beneath my fingers. The one with a horizontal crack. The
round hole that looked like a cigarette burn but almost certainly
wasn't. Cold keys on a winter's day. The subtle difference of texture
between the white and the black. Some tunes remain alive in my
head. Chord shapes and runs of notes live on in tactile memory.

The piano is long gone, but when I type, I remember the lessons
learned from other keys, writing with more passion than precision,
words that my fingers find as though they were still trying to release
the clear notes of long dead composers.
